\section{Реализация на гаранцията при седем хранилища}
\label{sec:enforcement}

Моделът на заплахите от глава~\ref{sec:threat-model} изисква гаранцията да се
провери поотделно за всяко хранилище, защото предпоставка 2 от
точка~\ref{sec:theory-compile-time} се изпълнява от конкретния конектор, а не от
общата абстракция. Настоящата глава описва накратко платформата, установява
откъде идва всяка част на заявката и проследява за всеки от седемте конектора как
стойността достига драйвера. Номерата на редовете се отнасят към версия
\code{a17f4e7} на хранилището с изходния код.

\subsection{Платформата}
\label{sec:enf-platform}

Изследваната библиотека е разработена в предходна дипломна работа за ОКС
„Бакалавър“~\cite{tsvetanov2026orm}, в която са описани обектният модел,
междинното представяне на заявките, асинхронното изпълнение и конекторите.
Настоящата работа не преповтаря тези резултати и не претендира за тях. Тя
използва библиотеката като платформа и изследва въпрос, който предходната работа
не поставя: доколко конструирането по време на компилация предотвратява
инжекциите и къде спира тази защита. Приносите, формулирани в заключението, се
отнасят само към това изследване и към промените в библиотеката, произтекли от
него.

За изложението са необходими три понятия от платформата. \textbf{Свойство}
(\code{orm::property<T, "name">}) описва колона: типът на стойността и името на
колоната са аргументи на шаблона. \textbf{Междинното представяне} на заявката
(\textit{IR}) е стойност от тип като \code{select\_query<...>}, чиито параметри на
шаблона кодират избраните колони, условията и ограниченията. \textbf{Конекторът}
превежда междинното представяне към езика на конкретно хранилище и предава
стойностите на драйвера.

\subsection{Откъде идва всяка част на заявката}
\label{sec:enf-origin}

Листинг~\ref{lst:enf-query} показва заявка, написана по документирания начин.

\begin{lstlisting}[language=C++,caption={Заявка, написана по документирания начин},label={lst:enf-query}]
struct Item
{
    orm::property<int, "id">           id;    // името на колоната е аргумент на шаблон
    orm::property<std::string, "name"> name;
};

namespace orm
{
    template <>
    struct table_name_trait<Item>
    {
        static constexpr std::string_view value = "items";   // константен израз
    };
}

constexpr auto q = orm::select(orm::field<&Item::id>)
                       .where(orm::field<&Item::name> == orm::Placeholder<std::string>{});

auto rows = db.execute(q, input);   // стойността пътува отделно от заявката
\end{lstlisting}

В листинга трябва да се забележат три неща: името на колоната е част от типа на
свойството, името на таблицата е константен израз, а стойността \code{input} не
участва в изграждането на \code{q} и се подава едва при изпълнението.
Таблица~\ref{tab:enf-origin} обобщава произхода на всяка част.

\begin{table}[H]
\centering
\caption{Произход на частите на заявката}
\label{tab:enf-origin}
\begin{tabularx}{\textwidth}{L{3.6cm}L{4.6cm}Y}
\toprule
\textbf{Част на заявката} & \textbf{Произход} & \textbf{Достъпна ли е за вход по време на изпълнение} \\
\midrule
Имена на таблици и колони & \code{constexpr} метаданни на указател към член (\code{details/member\_pointer.hpp:44}, \code{entity/property.hpp:19}) & Не \\
Форма на заявката (избор, условие, съединение, подреждане) & типове в междинното представяне & Не \\
Стойности & типизирани заместители \code{Placeholder<T>} и \code{ph<T, \_N>} & Да, само като данни \\
\bottomrule
\end{tabularx}
\end{table}

Таблицата съдържа твърдение, което не следва автоматично от синтаксиса на
листинга: че няма \textit{друг} път, по който стойност да достигне заявката. Това
твърдение се оказа невярно за предходната версия на библиотеката и е обект на
казуса в глава~\ref{sec:case-study}.

\subsection{Механизми при отделните хранилища}
\label{sec:enf-backends}

Гаранцията е една и съща за всички хранилища, но механизмът, по който тя се
постига, не е. Всеки конектор използва най-силното средство, което драйверът му
предоставя. Таблица~\ref{tab:enf-mechanisms} показва резултата от анализа на
изходния код.

\begin{table}[H]
\centering
\caption{Как стойността достига драйвера при всеки конектор}
\label{tab:enf-mechanisms}
\begin{tabularx}{\textwidth}{L{2.1cm}L{4.0cm}L{2.9cm}Y}
\toprule
\textbf{Хранилище} & \textbf{Механизъм} & \textbf{Вид} & \textbf{Място в кода} \\
\midrule
PostgreSQL & \code{PQexecParams} с масиви от параметри & свързване в сървъра & \code{postgresql\_live.hpp:552, 615} \\
MySQL & подготвени заявки с \code{MYSQL\_BIND} & свързване в сървъра & \code{mysql\_live.hpp:316} \\
SQLite & \code{sqlite3\_prepare\_v2} и \code{sqlite3\_bind\_*} & свързване & \code{sqlite\_db.hpp:62, 67, 72} \\
Cassandra & \code{cass\_statement\_bind\_*} & свързване в сървъра & \code{cassandra\_live.hpp:200-221} \\
Neo4j & параметри на Bolt чрез \code{neo4j\_map\_entry} & типизирани параметри на драйвера & \code{neo4j\_live.hpp:174} \\
Redis & форми с форматиращ низ и \code{argv} на \code{hiredis} & кодиране на протоколно ниво & \code{redis\_live.hpp:109, 180, 257} \\
MongoDB & построител на документи \code{bsoncxx} & структурно, без текстов език & \code{mongodb\_live.hpp:214-238} \\
\bottomrule
\end{tabularx}
\end{table}

От таблицата се извеждат четири вида механизъм. Четири конектора използват
свързване на параметри. Neo4j използва типизирани параметри на драйвера. Redis
разчита на кодирането на протокола. MongoDB изобщо не изгражда текст. Изводът е,
че една и съща абстракция се защитава пред четири различни протокола и
аргументът за сигурност трябва да е валиден поотделно за всеки от тях.

\subsubsection{Хранилища без текстов език за заявки}

При MongoDB и Redis стойността никога не става част от текст, който хранилището
анализира синтактично. MongoDB получава документ BSON, в който стойността е
поле със собствен тип. Redis получава кадри RESP, в които всеки аргумент е
предшестван от дължината си. Липсва граматика, от която низ на нападателя би
могъл да излезе. Тази позиция е по-силна от свързването на параметри, а не
равностойна на него, тъй като не зависи от правилно разделяне на маркери и
стойности. Разликата обикновено се губи, когато инжекциите в NoSQL се разглеждат
като допълнение към SQL инжекциите.

Тази позиция е валидна само докато стойността не се интерпретира като оператор.
При MongoDB конекторът добавя стойността като типизирано поле на документа
(\code{mongodb\_live.hpp:214-238}), а не разчита произволен документ от входа.
Затова низ, който съдържа обект с оператора {\ttfamily \$ne}, остава низ.

\subsubsection{Neo4j: форматиране, което не е интерполация}

Конекторът за Neo4j изгражда текста на Cypher чрез \code{std::format}, например
шаблона \code{MATCH (n:\{\}) RETURN \{\}\{\}} (\code{neo4j\_live.hpp:261}). Това
изглежда като сглобяване на низ и заслужава внимание. Вмъкваните полета обаче са
етикетът на възела и проекцията, и двете идват от \code{constexpr} метаданни на
типа, а не от входа. Стойностите по време на изпълнение се предават отделно,
като типизирани параметри на Bolt (\code{neo4j\_live.hpp:174}). Форматирането на
константи от времето на компилация в шаблон не е интерполация на недоверени
данни. Различието е съществено, но то трябва да се демонстрира, а не да се
приема. Докъде е стигнала тази демонстрация е разгледано в
точка~\ref{sec:eval-servers}.

\subsubsection{Redis: форматиращият низ не е \texttt{printf}}

Конекторът за Redis извиква \code{redisCommand} с форматиращ низ, в който
аргументът е означен с {\ttfamily \%s} (\code{redis\_live.hpp:180}). Въпреки
сходството с \code{printf}, \code{hiredis} не замества текст. Всеки аргумент се
кодира като отделен елемент на RESP с префикс за дължина, така че съдържанието му
не може да бъде прочетено като протокол. Формата с \code{argv}
(\code{redis\_live.hpp:257}) прави това изрично.

\subsubsection{Управление на транзакции}

Командите за транзакции изпращат литералните константи \code{BEGIN},
\code{COMMIT} и \code{ROLLBACK} (\code{postgresql\_live.hpp:506, 518, 530}). Те
не съдържат потребителски данни и не нарушават гаранцията.

\subsection{Проверка на границата}
\label{sec:enf-boundary}

Анализът по хранилища потвърждава предпоставка 2 за всеки конектор поотделно, но
оставя един въпрос: стига ли всеки път към конектор само до добре формирано
междинно представяне. В хода на настоящата работа беше установено, че отговорът е
бил отрицателен (глава~\ref{sec:case-study}). Затова в библиотеката е добавена
проверка, извиквана от всяка входна точка, която достига конектор.

\begin{lstlisting}[language=C++,caption={Проверката на границата (съкратено)},label={lst:enf-shape}]
template <typename Query>
constexpr void check_query_shape()
{
    using Q = std::remove_cvref_t<Query>;
    static_assert(is_query_ir<Q>::value,
        "orm: this is not a query. ...");
    if constexpr (is_query_ir<Q>::value)
        static_assert(!query_has_literal<Q>::value,
            "orm: a literal operand reaches execution. ...");
}
\end{lstlisting}

Проверката в листинг~\ref{lst:enf-shape} има две части. Първата отхвърля всичко,
което не е един от четирите вида заявки (\code{select\_query},
\code{insert\_query}, \code{update\_query}, \code{delete\_query}). Анализът на
всички конектори показа, че това са единствените видове, които те изпълняват,
така че проверката не отхвърля нито една легитимна употреба. Втората част
обхожда условията, съединенията и присвояванията в заявката и отхвърля всеки
\textbf{литерален операнд}: стойност, записана направо в заявката, вместо през
заместител. Сравнението с \code{nullptr} и сравнението между две колони не са
литерали и остават разрешени.

Проверката се извиква от седемте входни точки, които достигат конектор:
\code{db::operator<<}, \code{db::execute}, \code{db::find\_one},
\code{db::prepare}, \code{prepared\_query::execute} и двете асинхронни
съответствия в \code{async\_db}. Пропускането на която и да е от тях би оставило
заобиколен път. Затова всяка входна точка се покрива от отделен компилационен
тест (глава~\ref{sec:evaluation}).

\subsection{Известен остатъчен риск}
\label{sec:enf-residual}

Анализът установи и един път, който е по-слаб от свързването, без да е
използван. Функцията \code{to\_sql\_literal} (\code{mysql\_live.hpp:249})
изгражда SQL чрез \textbf{екраниране} и вмъкване на стойността, вместо чрез
свързване. Коментарът ѝ я описва като предназначена за заявки без подготвени
изрази. Търсене в каталозите \code{lib} и \code{tests} намира само дефиницията и
нито едно извикване. Следователно тя е мъртъв код и не засяга текущата гаранция.
Тя е отбелязана, защото екранирането е по-слаб механизъм
(точка~\ref{sec:theory-remedies}) и бъдещо използване би отслабило гаранцията за
MySQL, без това да е видимо от интерфейса. Препоръката е функцията да бъде
премахната или изрично изолирана.
